% chapters/10/evaluasi.tex
\bagian{Asesmen dan Instrumen Penilaian}
\label{sec:asesmen-astar}

\begin{asesmenbox}
\textbf{Kedudukan asesmen bab ini.} Bab ini menuntaskan \SUBCPMK{2.3}.
Ketiga indikatornya diperiksa melalui penilaian artefak program (Tugas
Praktikum 08), tangkapan layar pencarian, dan soal konsep yang menuntut
analisis.

\vspace{0.4em}
Materi Bab~X termasuk bahan \textbf{Mini Proyek AI} pada minggu kesembilan dan
bahan \textbf{Ujian Tengah Semester} pada minggu kesepuluh.
\end{asesmenbox}

\subsection*{Instrumen 1: Rubrik Praktikum 08}

\begin{table}[htbp]
\centering
\caption{Rubrik penilaian Praktikum 08 beserta indikator yang diukur}
\label{tab:rubrik-praktikum-08}
{\setlength{\tabcolsep}{4pt}\footnotesize%
\begin{tabularx}{\textwidth}{@{}L{3.2cm} C{1.3cm} C{2.1cm} Y@{}}
\toprule
\headrow \textbf{Kriteria} & \textbf{Bobot} & \textbf{Indikator} & \textbf{Deskripsi capaian} \\
\midrule
Representasi peta dan pencarian & 35\% & \ind{a} &
Peta dibaca menjadi graf dengan benar, fungsi tetangga menyaring batas dan
dinding, dan pencarian A* memakai antrean prioritas. \\
\addlinespace
Analisis pengaruh heuristik & 35\% & \ind{b} &
Tabel perbandingan tiga heuristik terisi lengkap, biaya jalur terbukti sama,
dan pengaruh heuristik terhadap kinerja dijelaskan beserta akibat heuristik
yang melebihkan jarak. \\
\addlinespace
Rekonstruksi dan visualisasi & 30\% & \ind{c} &
Jalur direkonstruksi dengan benar, pemeriksaan keabsahan berjalan bersih, dan
tangkapan layar pencarian disertakan pada laporan. \\
\bottomrule
\end{tabularx}}
\end{table}

\subsection*{Instrumen 2: Soal Analisis}

Jawablah dengan singkat. Kunci jawaban beserta indikator yang diukur disajikan
pada Tabel~\ref{tab:kunci-konsep-ch10}.

\begin{enumerate}
  \item Uraikan bagaimana sebuah peta berpetak menjadi graf, dan hitunglah
        banyaknya tetangga yang mungkin pada satu petak untuk gerak empat arah.
  \item Sebuah petak berjarak delapan langkah dari titik awal dan berjarak tiga
        satuan Manhattan dari titik tujuan. Hitunglah $f$-nya, lalu jelaskan
        mengapa nilai itu menentukan urutan pemeriksaannya.
  \item Pada peta bab ini, Manhattan memeriksa 274 petak sedangkan Dijkstra 437
        petak, tetapi biaya jalurnya sama. Jelaskan mengapa demikian.
  \item Sebuah heuristik dikalikan dua agar pencariannya lebih cepat. Jelaskan
        akibatnya terhadap mutu jalur, dan bagaimana Anda mengenalinya.
  \item Tuliskan langkah-langkah rekonstruksi jalur dari tabel asal, termasuk
        langkah yang diperlukan apabila tujuan tidak dapat dicapai.
\end{enumerate}

\begin{table}[htbp]
\centering
\caption{Kunci soal analisis Bab X beserta indikator yang diukur}
\label{tab:kunci-konsep-ch10}
{\setlength{\tabcolsep}{4pt}\footnotesize%
\begin{tabularx}{\textwidth}{@{}C{0.8cm} C{2.1cm} Y@{}}
\toprule
\headrow \textbf{No} & \textbf{Indikator} & \textbf{Pokok jawaban} \\
\midrule
1 & \ind{a} &
Setiap petak kosong menjadi simpul, setiap pasangan petak bersebelahan menjadi
sisi berbiaya satu, dan dinding tidak menjadi simpul. Untuk gerak empat arah,
setiap petak memiliki paling banyak empat tetangga. \\
\addlinespace
2 & \ind{b} &
$f = g + h = 8 + 3 = 11$. Petak dengan $f$ terkecil diperiksa lebih dahulu,
sebab nilai itu menyatakan perkiraan biaya total jalur yang melalui petak
tersebut. \\
\addlinespace
3 & \ind{b} &
Kedua heuristik itu tidak melebihkan jarak, sehingga keduanya menghasilkan
jalur terpendek dengan biaya yang sama. Yang berbeda hanyalah banyaknya petak
yang diperiksa, sebab Manhattan lebih mengarahkan pencarian ke arah tujuan
daripada tanpa perkiraan. \\
\addlinespace
4 & \ind{b} &
Heuristik yang melebihkan jarak tidak lagi memenuhi syarat, sehingga jalur yang
dihasilkan dapat lebih panjang daripada yang terpendek. Pencariannya memang
lebih cepat, dan tanda satu-satunya adalah biaya jalur yang lebih besar pada
tabel perbandingan. \\
\addlinespace
5 & \ind{c} &
Periksa lebih dahulu apakah tujuan ada pada tabel asal; apabila tidak, jalurnya
kosong dan keadaan itu harus ditangani. Apabila ada, telusuri dari tujuan ke
asalnya berulang kali sampai asalnya bernilai kosong, kumpulkan petaknya, lalu
balik urutannya. \\
\bottomrule
\end{tabularx}}
\end{table}
